% Generated draft: requires mathematical and source review.
\documentclass[11pt]{article}
\usepackage[margin=25mm]{geometry}
\usepackage{fontspec}
\setmainfont{Latin Modern Roman}
\usepackage{amsmath,amssymb,mathtools,unicode-math}
\setmathfont{Latin Modern Math}
\usepackage{enumitem,needspace,xurl,hyperref}
\setlength{\parindent}{0pt}
\setlength{\parskip}{5pt}
\setlength{\emergencystretch}{3em}
\newcommand{\N}{\mathbb{N}}
\newcommand{\Z}{\mathbb{Z}}
\newcommand{\Q}{\mathbb{Q}}
\newcommand{\R}{\mathbb{R}}
\newcommand{\C}{\mathbb{C}}
\newcommand{\F}{\mathbb{F}}
\newcommand{\A}{\mathbb{A}}
\newcommand{\PP}{\mathbb{P}}
\newcommand{\E}{\mathbb{E}}
\newcommand{\eps}{\varepsilon}
\newcommand{\dd}{\,\mathrm{d}}
\begin{document}
\section*{UnsolvedMath ID: 800526}
\textbf{Source problem code:} catalogue entry
\subsection{Definitions and mathematical statement}
Let \(A\in\R^{n\times d}\), \(b\in\R^n\), and \(\Omega\in\R^{n\times k}\), where \(n,d,k\) are positive integers. The source does not specify the range of \(p\). Under the usual norm convention, for \(v\in\R^m\),
\[
\|v\|_{\ell_p}=\left(\sum_{i=1}^m |v_i|^p\right)^{1/p}\quad(1\le p<\infty),
\qquad
\|v\|_{\ell_\infty}=\max_{1\le i\le m}|v_i|.
\]
These definitions do not settle which values of \(p\), including whether \(p=\infty\), the source intends to include.

Define the optimal original residual and the set of sketch-and-solve solutions by
\[
r_p(A,b)=\min_{x\in\R^d}\|Ax-b\|_{\ell_p},
\qquad
S_p(A,b,\Omega)=\operatorname*{argmin}_{x\in\R^d}
\|\Omega^\top Ax-\Omega^\top b\|_{\ell_p}.
\]
The original residual lies in \(\R^n\), and the sketched residual lies in \(\R^k\).

\textbf{Problem.} Formulate a generalization of the source's OSI property to \(\ell_p\) regression that ensures constant-factor accuracy of the sketch-and-solve procedure: for a solution \(\tilde x\in S_p(A,b,\Omega)\), the proposed property should imply
\[
\|A\tilde x-b\|_{\ell_p}\le C\,r_p(A,b)
\]
for a constant \(C\ge1\). In particular, when \(r_p(A,b)=0\), this inequality requires \(A\tilde x=b\).

The supplied record does not define OSI, specify the allowed dependence of \(C\), or state whether the guarantee must hold for every sketched minimizer or for a specified selection rule. It also gives no distribution for \(\Omega\), success probability, sketch-size requirement, or quantifier order governing whether a sketch must work for a fixed instance or uniformly over a class of instances. These aspects of the requested property remain unresolved; no particular generalization is asserted here.
\subsection{Short English statement}
Find an \(\ell_p\) analogue of the source's OSI property that guarantees that solving the sketched regression problem produces a solution whose original residual is within a constant factor of optimal. The supplied record leaves OSI, the intended range of \(p\), and several guarantee quantifiers unspecified.
\subsection{Sources}
\begin{itemize}
\item Local UnsolvedMath registry, permanent ID \(800526\), source problem code \texttt{catalogue entry}, titled ``Linear Systems and Eigenvalue Problems --- Problem 35 (Regression in the \(\ell_p\) norm).'' The registry attributes the problem to Amsel, Baumann et al., with proposed year 2026, and labels it open. This is a supplied status label, not an independently verified mathematical or literature assessment.
\item Amsel, Baumann et al., \emph{Linear Systems and Eigenvalue Problems: Open Questions from a Simons Workshop} (2026), Problem 35. Supplied source URL: \url{https://arxiv.org/abs/2602.05394}. Supplied attribution/context: ``Regression in the \(\ell_p\) norm.''
\item Supplied provenance: accessed 2026-10-16; extraction method \texttt{arxiv\_latex}; upstream revision \texttt{d371ee917ac9647fbb06cd6fc529deec6f3dda82}. These are supplied source-verification metadata; the reference and its contents were not independently checked for this draft.
\end{itemize}
\end{document}
